% Introduction generale non numerotee
\phantomsection
\markboth{Introduction générale}{Introduction générale}
\addcontentsline{toc}{chapter}{Introduction générale}
\begin{center}
{\Large\bfseries Introduction générale}
\end{center}
\vspace{0.4cm}

L'\textbf{Internet des Objets} (IoT, \emph{Internet of Things}) désigne l'interconnexion de dispositifs physiques capables de collecter, transmettre et exploiter des données dans des environnements très variés : villes intelligentes, santé connectée, industrie, agriculture, transport ou infrastructures de télécommunications~\cite{meneghello2019iot}. Ces dispositifs peuvent être de simples capteurs, des passerelles industrielles, des compteurs intelligents ou des équipements embarqués. Leur rôle est de rapprocher le monde physique des systèmes numériques afin d'automatiser la surveillance, l'analyse et la prise de décision.

Cette évolution apporte cependant de nouveaux risques. Contrairement aux systèmes informatiques classiques, les objets connectés sont souvent nombreux, hétérogènes, parfois peu puissants et rarement administrés manuellement après leur déploiement. Ils communiquent sur des réseaux partagés, publient des données sensibles et s'appuient sur des protocoles applicatifs légers comme le \textbf{Message Queuing Telemetry Transport} (MQTT) ou le \textbf{Constrained Application Protocol} (CoAP). La moindre faiblesse d'authentification, de chiffrement ou de segmentation peut donc exposer l'ensemble de l'écosystème à l'interception, à l'usurpation d'identité, au déni de service ou à l'exfiltration de données.

La sécurité des communications IoT doit ainsi être pensée de bout en bout. Elle ne se limite pas à chiffrer un canal réseau : elle implique la gestion fiable des identités, l'émission et la rotation de certificats, l'authentification mutuelle entre dispositifs et services, la restriction des accès par politiques, ainsi qu'une supervision capable de détecter les comportements anormaux. Dans un contexte opérateur, ces exigences deviennent encore plus importantes, car l'infrastructure peut héberger plusieurs services, plusieurs clients et des volumes élevés de messages.

Ce projet de fin d'études s'attaque à un problème précis et récurrent dans les déploiements IoT opérateurs~: l'absence d'une chaîne de sécurité cohérente, de bout en bout, allant de l'identité cryptographique de chaque dispositif jusqu'à la détection automatique de comportements anormaux. Les solutions existantes couvrent généralement un seul maillon~--- chiffrement \emph{ou} supervision \emph{ou} détection~--- sans les intégrer ni les valider conjointement dans un environnement représentatif. Il en résulte des angles morts~: un dispositif portant un certificat valide mais volé peut se connecter sans être détecté~; un flux chiffré peut transporter une exfiltration invisible aux règles statiques~; une alerte IDS isolée ne suffit pas à qualifier une attaque sans corrélation SIEM. Ce projet vise à combler ces lacunes en construisant et en validant expérimentalement une architecture de sécurisation intégrée, reproductible et mesurable, ancrée dans les contraintes concrètes d'un opérateur télécom tel que Tunisie Telecom.


Le rapport est organisé en cinq chapitres. Le chapitre~\ref{chap:contexte-general} présente le contexte général du projet, l'organisme d'accueil, la problématique, les objectifs et la solution proposée. Le chapitre~\ref{chap:etat-art} expose l'état de l'art relatif aux menaces IoT, aux protocoles, aux mécanismes de sécurité et aux outils de supervision. Le chapitre~\ref{chap:conception} détaille la conception : besoins, choix méthodologique, planification Scrum, backlog et diagrammes en \textbf{Unified Modeling Language} (UML). Le chapitre~\ref{chap:realisation} décrit la réalisation technique de l'architecture, les scénarios d'attaque, le pipeline d'\textbf{apprentissage automatique} (ML, \emph{Machine Learning}) et les captures d'interfaces. Enfin, le chapitre~\ref{chap:tests-conclusion} présente les tests, les résultats, les limites et les perspectives. Une conclusion générale synthétise ensuite les apports du projet.
